\subsection{The output of the application, panel by panel}
The layer \emph{Soil, plant and air} is the picture of the whole chain that Chapter 8 describes in words: one continuous water column from the groundwater to the atmosphere, in which the suction rises from $-0.02$ MPa in the soil to $-56$ MPa in the air.

Each height of the column from the water table to the leaves has its own suction and content; the arrows show the direction of the flow.

\appfig{spac_0}{0.96}{Soil to sky: one continuous water column, with the suction, the water content and the pores holding water at each height.}

The figure compares the pull of the atmosphere on the water with the hold of the soil on it and shows that the water moves along the gradient of the total head.

\appfig{spac_1}{0.96}{The air pulls with $-56$ MPa; the soil holds on with $-0.02$ MPa.}

The three profiles are the simulated state of the column on the day of the replay.

\appfig{spac_2}{0.96}{Simulated suction, water content and pores holding water along the column.}

